%=========== ΛΥΣΗ - 43ο ΘΕΜΑ Α ===========
\providecommand{\theoriaref}[3]{%
\begin{center}
\fcolorbox{maincolor!70!black}{maincolor!8}{%
\begin{minipage}{.88\linewidth}
\textcolor{maincolor!80!black}{\kerkissans{\textbf{#1~\ref{#2}}}}\\[1mm]
#3
\end{minipage}}
\end{center}}
\providecommand{\orismosref}[1]{%
\theoriaref{Ορισμός}{#1}{Η απάντηση δίνεται από τον αντίστοιχο ορισμό.}}
\providecommand{\apodeixiref}[1]{%
\theoriaref{Θεώρημα}{#1}{Η ζητούμενη απόδειξη δίνεται στην αντίστοιχη απόδειξη.}}

\begin{thema}{Α}
\begin{erwthma}
\item \apodeixiref{thewrhma:themelivdes_oloklhrwtikou}

\item \orismosref{orismos:synthesh_synarthsewn}

\item \orismosref{orismos:katakoryfh_asymptwth}

\item Οι χαρακτηρισμοί των προτάσεων είναι:
\begin{alist}
\item \textbf{Λάθος}. Η μορφή $0\cdot(+\infty)$ είναι απροσδιόριστη και δεν οδηγεί πάντοτε στο $0$.
\item \textbf{Λάθος}. Το θεώρημα \eng{Rolle} είναι ειδική περίπτωση του Θ.Μ.Τ. όταν $f(a)=f(\beta)$, όχι το αντίστροφο.
\item \textbf{Σωστό}. Αν $f(x)\geq g(x)$ στο $[a,\beta]$, τότε το εμβαδόν ανάμεσα στις δύο καμπύλες είναι
\[ E=\dint_a^\beta (f(x)-g(x))\d x. \]
\item \textbf{Σωστό}. Όταν η $f$ είναι κοίλη σε ένα διάστημα, η παράγωγός της $f'$ είναι γνησίως φθίνουσα σε αυτό.
\item \textbf{Σωστό}. Για παράδειγμα, η $f(x)=|x|$ είναι συνεχής στο $\mathbb{R}$, αλλά στο $0$ δεν έχει εφαπτομένη με ορισμένο συντελεστή διεύθυνσης.
\end{alist}

\item Σωστή απάντηση είναι η \textbf{β}. Θέτουμε
\[ u=\hm x,\qquad \d u=\syn x\,\d x, \]
οπότε το ολοκλήρωμα αντιμετωπίζεται με αντικατάσταση $u=\hm x$.
\end{erwthma}
\end{thema}
%=========== ΤΕΛΟΣ ΛΥΣΗΣ - 43ο ΘΕΜΑ Α ===========
